\documentclass[11pt]{article}
\usepackage[a4paper,margin=25mm]{geometry}
\usepackage[T1]{fontenc}
\usepackage{lmodern,amsmath,amssymb,amsthm,booktabs,microtype}
\usepackage[hidelinks]{hyperref}
\hypersetup{pdftitle={Exact Key-List Recovery in SEPAR},pdfauthor={Noah Ostle}}
\newtheorem{theorem}{Theorem}[section]
\newtheorem{lemma}[theorem]{Lemma}
\newtheorem{proposition}[theorem]{Proposition}
\newcommand{\E}{\mathcal E}
\newcommand{\D}{\mathcal D}
\newcommand{\xor}{\mathbin\oplus}
\newcommand{\OR}{\mathbin{\mathrm{OR}}}
\newcommand{\rol}{\operatorname{rol}_{16}}
\newcommand{\pack}{\operatorname{pack}}
\newcommand{\calP}{\mathcal P}
\newcommand{\Tr}{\operatorname T}
\title{Exact Key-List Recovery in SEPAR}
\author{Noah Ostle}
\date{19 October 2026}
\begin{document}
\maketitle
\begin{abstract}
SEPAR's word permutations preserve nested blocks. State additions
disturb these blocks only at narrow carry boundaries. We use this
structure to recover numerical rotor keys from quotient paths, then
remove rotors inward by an exact partition join. Initialization
supplies the remaining state words and the first two rotor keys.
The result is a chosen-IV algorithm that returns a list of complete
256-bit keys without ranking or candidate cutoffs. For every fixed
master key, its true key is in the list with probability at least
\(1-1.307164\cdot10^{-8}\). We also give a 511-query
distinguisher.
\end{abstract}

\section{Target and attack model}
We analyze the literal supplied C implementation of SEPAR~\cite{separ}.
It has eight keyed word permutations, a 128-bit state, and a 16-bit
LFSR. The published description and the C source differ in some
sequential assignments; the equations here follow the source.
An adversary may choose an IV and a plaintext, reset, and ask for the
first ciphertext word repeatedly under that IV. This reuse is part of
the attack model. The recovery never reads the internal state or a
rotor oracle. Complete candidate keys are also checked against saved
public ciphertext tables.

All additions below are modulo \(2^{16}\). Write a word as four
nibbles \((a,b,c,d)\), from most to least significant. Let
\(\Sigma\) apply the four S-boxes in parallel. Their entries for
inputs \(0,\ldots,F\) are
\[
\begin{array}{c|*{16}{c}}
S_1&1&F&B&2&0&3&5&8&6&9&C&7&D&A&E&4\\
S_2&6&A&F&4&E&D&9&2&1&7&C&B&0&3&5&8\\
S_3&C&2&6&1&0&3&5&8&7&9&B&E&A&D&F&4\\
S_4&D&B&2&7&0&3&5&8&6&C&F&1&A&4&9&E
\end{array}
\]
Rotor \(i\) has its own two-word key \(K_i=(k_{i,0},k_{i,1})\).
Its four subkeys are
\[
\begin{split}
G(u)&=u\OR\bigl(S_1((u\gg6)\mathbin{\&}15)\ll6\bigr),\\
\kappa_0&=k_{i,0},\qquad \kappa_1=k_{i,1},\\
\kappa_2&=G(\rol(k_{i,0},6))\xor(i+2),\\
\kappa_3&=G(\rol(k_{i,1},10))\xor(i+3).
\end{split}
\]
For \(j=0,1,2,3\), apply \(x\leftarrow L\Sigma(x\xor\kappa_j)\);
finally apply
\(x\leftarrow\Sigma(x\xor\kappa_0\xor\kappa_1)
\xor\kappa_2\xor\kappa_3\).
The implemented linear map is
\begin{equation}
                 L(a,b,c,d)=(b\xor c,b,a,a\xor d).
\label{eq:linear}
\end{equation}
This defines a permutation \(\E_i\), with inverse \(\D_i\).
The OR in \(G\) is intentional: replacing bits instead would
describe a different target.

For an initialized state \(s=(s_1,\ldots,s_8)\), the first-word map is
\begin{equation}
 R_0(p)=p,\qquad R_i(p)=\E_i(R_{i-1}(p)+s_i),\qquad C(p)=R_8(p).
 \label{eq:cascade}
\end{equation}
Initialization consists of four rounds of the map given in
Section~\ref{sec:state}. Later words use a state and LFSR update
specified by the source. All acquisition queries in this attack
reset before the first word, so \eqref{eq:cascade} is the map under
analysis. Complete-key replay uses the source's initialization.

\section{Preserved blocks and useful states}\label{sec:state}
\begin{lemma}[Triangular rotor]\label{lem:triangular}
Every numerical rotor key gives
\[
 \E_i(a,b,c,d)
 =\bigl(P_{i,b}(a),\,Q_i(b),\,M_{i,b}(c),\,
         T_{i,a,b,c}(d)\bigr),
\]
where each one-coordinate map is a permutation.
Consequently, the high twelve output bits depend only on the high
twelve input bits. The middle-byte quotient
\[
 m_i(16b+c)=16Q_i(b)+M_{i,b}(c)
\]
is a permutation of 256 values.
\end{lemma}
\begin{proof}
One keyed substitution and linear layer send
\((a,b,c,d)\) to coordinates depending respectively on
\((b,c),b,a,(a,d)\). Four layers give dependency sets
\((a,b),b,(b,c),(a,b,c,d)\). The final substitution and XOR
do not enlarge them. With the other coordinates fixed, each
coordinate path is a composition of nibble permutations and XORs.
\end{proof}
In particular, a change confined to \(d\) stays in \(d\) under
\(L\); its nibble branch number is two. A carry at a state addition
may move an input between two adjacent quotient blocks, but cannot
diffuse it through all blocks. These narrow carry boundaries are the
objects decoded below. The OR schedule also leaves small raw-key
fibres after quotient data have been fixed.

One initialization round first computes
\[
v_1=\E_1(s_1+s_3+s_5+s_7),\qquad
v_i=\E_i(v_{i-1}+s_i)\quad(2\le i\le8),
\]
then replaces \(s_1\) by \(s_1+v_8\) and \(s_i\) by
\(s_i+v_{i-1}\) for \(i\ge2\).
\begin{lemma}[Uniform state]\label{lem:uniform}
For every fixed key, four initialization rounds permute the
128-bit IV space. Independent uniform IVs therefore give
independent uniform initialized state words.
\end{lemma}
\begin{proof}
Given the output \(t\) of one round, invert in the order
\[
 s_1=t_1-\E_8(t_8),\quad
 s_i=t_i-\E_{i-1}(t_{i-1})\ (i=8,\ldots,3),\quad
 s_2=t_2-\E_1(s_1+s_3+s_5+s_7).
\]
Four bijections compose to a bijection.
\end{proof}

Let a block have \(B\) consecutive words. An input row consists
of all its \(B\) words. The next lemma is the same induction at
every quotient width used by the attack.
\begin{lemma}[Perturbed blocks]\label{lem:block}
Align a row with the first state addition. Write the intervening
state additions before a chosen last rotor as \(Bh_i+e_i\), with
signed residues \(e_i\), and put \(S=\sum_i|e_i|\).
If the final state residue is \(b\), with
\begin{equation}
              S<\min(b,B-b),
\label{eq:support}
\end{equation}
the two endpoints of the last quotient-cycle edge both occur in
the row, with at least \(B-S\) points between them. At least
\(B-S\) inverse images under the last rotor lie in one interval
\([Bv+b,Bv+b+B-1]\). If
\(2S<\min(b,B-b)\), the endpoints are the unique two largest
quotient-label counts.
\end{lemma}
\begin{proof}
After the first rotor, an aligned row fills one block. At the
next translation, at most \(|e_i|\) distinct low coordinates cross
a boundary. Discard them and apply the rotor; it permutes the
remaining coordinates within a nominal block. Induction leaves
at least \(B-S\) points, and the nominal block labels are a
permutation of their original labels. The final addition by \(b\)
splits the full nominal block into parts of sizes \(b,B-b\).
Condition~\eqref{eq:support} leaves each part nonempty after the
discarded points are removed. All retained points lie in the
stated translated interval before the final rotor. If \(2S\)
is below each part size, each intended label has more points
than all the other labels combined, proving dominance.
\end{proof}

\section{A preliminary distinguisher}
The block weakness already distinguishes SEPAR without numerical
key recovery.
\begin{theorem}[511-query distinguisher]\label{thm:dist}
At one IV, query \(p=0,\ldots,510\), resetting before each query.
Accept if some consecutive 256-input window has at most 128
ciphertext high-byte values. Every SEPAR key and IV is accepted.
A uniform random word permutation is accepted with probability
at most \(1.762252054731\cdot10^{-9}<2^{-29}\).
\end{theorem}
\begin{proof}
One of the 256 windows aligns \(p+s_1\) with a complete high-byte
block. Equation~\eqref{eq:linear} makes the first rotor's high byte
constant on it. Each of the seven later additions can split a
high-byte block into at most two, while a rotor merely permutes
the classes. Thus the ciphertext uses at most \(2^7\) high bytes.

For the ideal permutation, let \(A_{r,k}\) count ordered selections
of \(r\) distinct words occupying \(k\) high-byte classes. With
\(A_{0,0}=1\) and infeasible entries zero,
\[
 A_{r+1,k}=A_{r,k}(256k-r)
          +256(257-k)A_{r,k-1}.
\]
The two terms reuse a class or open a new one. Write
\((a)_b=a(a-1)\cdots(a-b+1)\). The probability for one window is
\[
              \frac{\sum_{k\le128}A_{256,k}}{(65536)_{256}}.
\]
A union bound over the 256 windows and exact integer evaluation of
this recurrence give the stated value.
\end{proof}
We now turn to complete key recovery.

\section{Decoding a rotor from middle paths}\label{sec:decoder}
This section supplies a numerical decoder shared by outer exposure
and the inward factorization. For a rotor subkey write
\(\kappa_j=(\alpha_j,\beta_j,\gamma_j,\delta_j)\). Define
\[
J^{u,v,w,t}(x)=
S_3\!\left(S_1\!\left(S_3\!\left(S_1(S_3(x)\xor u)\xor v\right)
\xor w\right)\xor t\right).
\]
Substitution into the four literal rounds yields
\begin{align}
M_b(c)&=(\gamma_2\xor\gamma_3)\xor
       J^{u_b,v,w_b,t}(c\xor\gamma_0),\label{eq:middlecore}\\
u_b&=S_2(b\xor\beta_0)\xor\alpha_1,\quad
v=\gamma_2,\quad t=\gamma_0\xor\gamma_1,\nonumber\\
w_b&=S_2(S_2(S_2(b\xor\beta_0)\xor\beta_1)\xor\beta_2)
       \xor\alpha_3.\nonumber
\end{align}
Thus, for each output lane \(y=Q(b)+\lambda\bmod16\), the
sequence \(M_b(0),\ldots,M_b(15)\) is a path through its sixteen
vertices. Its orientation and the common cyclic lane translation
\(\lambda\) may be unknown.

For a nibble table \(f\), put
\[
 \pack(f)=\sum_{x=0}^{15}f(x)2^{4x},\qquad
 \mathcal C(f)=
 \min_{a=0}^{15}\pack\bigl(x\mapsto f(x\xor a)\xor f(a)\bigr).
\]
This canonicalizes input and output XOR. Reversing a path means
\(c\mapsto15-c=c\xor15\), so it does not change the canonical value.
Let \(I=S_2^{-1}\) and
\begin{equation}
F_{j,r,s}(y)=I(I(I(I(y\xor r\xor s)\xor j)\xor s)\xor r).
\label{eq:laneinverse}
\end{equation}
Two finite identities, computed from the S-boxes printed above, are
used throughout:
\begin{center}
\begin{tabular}{lrr}
\toprule
Family & Tables & Distinct signatures\\
\midrule
\(\mathcal C(J^{u,v,w,t})\) & 65536 & 65536\\
\(\pack(y\mapsto F_{j,r,s}(y-\lambda)\xor
 F_{j,r,s}(-\lambda))\) & 65536 & 65536\\
\bottomrule
\end{tabular}
\end{center}
The first calculation loops over all four nibbles
\((u,v,w,t)\), evaluates all sixteen entries, and sorts their
canonical 64-bit words. The second loops over
\((j,r,s,\lambda)\), evaluates \eqref{eq:laneinverse} on all sixteen
inputs, and sorts those packed words. Equality of adjacent sorted
entries is the complete collision test. Seven first-family
tables have a nonidentity input/output XOR stabilizer. These are
exhaustive finite assertions, not estimates over master keys.

\begin{theorem}[Forty numerical keys]\label{thm:forty}
Sixteen supplied unoriented middle paths, indexed by output lane
up to a common cyclic translation, determine at most forty raw
32-bit rotor keys. The exact decoder emits all compatible keys and
the compatible translation. It includes the true key for true paths.
\end{theorem}
\begin{proof}
Canonicalize each path. The first finite identity determines its
\((u_y,v_y,w_y,t_y)\); reject if the \(v_y,t_y\) disagree.
The \(u_y\) must be distinct. Set
\(j=\beta_0\xor\beta_1\), \(r=\beta_2\),
\(s=\beta_3\), and \(\kappa=\beta_1\xor\alpha_1\).
Inverting the autonomous output lane gives
\[
                  u_y=F_{j,r,s}(y-\lambda)\xor\kappa .
\]
The second finite identity applied to \(u_y\xor u_0\) fixes
\((j,r,s,\lambda)\); \(u_0\) fixes \(\kappa\). Then
\[
 \alpha_3=w_y\xor S_2(S_2(u_y\xor\kappa)\xor r)
\]
must be the same in every lane.

At least one path core has no nonidentity stabilizer: there are
sixteen distinct \(u_y\) and only seven exceptional cores in the
entire family. This path fixes \(\gamma_0\) up to reversal;
the key schedule fixes \(\gamma_0\bmod4=r\gg2\), choosing one
orientation. Equation~\eqref{eq:middlecore} then fixes \(\gamma_3\).
The implementation checks the equations on every row, so it
does not select a special row or key value.

Write \(g(x)=x\OR S_1(x)\). The schedule gives
\[
 \beta_1\gg2=\gamma_3\mathbin{\&}3,\quad
 g(\alpha_1)=(\gamma_3\gg2)+4(s\mathbin{\&}3),\quad
 \alpha_1=\beta_1\xor\kappa .
\]
For each fixed upper two bits, every fibre of \(g\) has size at most
two, as is checked directly from
\[
(g(0),\ldots,g(15))=
(1,15,11,3,4,7,7,15,14,9,14,15,13,15,14,15).
\]
There are therefore at most two \(\beta_1\), four lower-two-bit
choices for \(\alpha_0\), and five choices for \(\delta_0\)
after its \(g\)-value is fixed. The remaining raw-key nibbles
follow from the displayed schedule. Checking the complete paths
leaves at most \(2\cdot4\cdot5=40\) keys, all of which are emitted.
\end{proof}

For a complete quotient cycle let
\(q_K(x)=\E_K(16x)\gg4\). A translated cycle has order
\begin{equation}
 q_K(\varepsilon x+a)+h\pmod {4096},
 \qquad \varepsilon\in\{1,-1\},
\label{eq:cycle}
\end{equation}
up to its cyclic starting point. Try \(h\bmod16\), reduce the
result modulo 256, and extract its sixteen middle paths.
Theorem~\ref{thm:forty} fixes \(h/16\bmod16\).
For each emitted key, try the sixteen remaining top-nibble
translations and compare all 4096 entries in both directions.
This finite procedure emits every compatible numerical rotor and
full quotient translation; it has no ranking cutoff.

\section{Recovering the outer rotor}\label{sec:outer}
We use every pair of quotient labels supported by a row.
This avoids a dominance assumption: false edges may enter the
graph, but the true path remains.

Align a row of sixteen plaintexts at \(r=-s_1\bmod16\).
For \(i=2,\ldots,7\), write \(s_i=16h_i+e_i\), using signed
residues, and let \(S=\sum_{i=2}^7|e_i|\).
Put \(b=s_8\bmod16\). The outer event is
\begin{equation}
                    S<\min(b,16-b).
\label{eq:outerevent}
\end{equation}
The condition itself forces \(S\le7\), so no separate norm cutoff
is needed. A signed residue with \(|e_i|\le7\) has a unique
representative modulo sixteen.

\begin{theorem}[Supported outer paths]\label{thm:outer}
On event~\eqref{eq:outerevent}, a public graph made from 256 aligned
16-point rows contains all fifteen internal edges of each of the
sixteen true middle paths of \(K_8\). Enumerating graph paths,
decoding them by Theorem~\ref{thm:forty}, and applying the common
carry test below retains the true numerical \(K_8\).
The event has exact probability
\begin{equation}
                     p_{\rm out}=\frac{48639}{16^7}
                     =0.00018119439482688904
\label{eq:pout}
\end{equation}
for every fixed master key.
\end{theorem}
\begin{proof}
Lemma~\ref{lem:block} says that each row contains two adjacent
outer quotient labels with combined count at least \(16-S\ge9\).
Project a top-twelve-bit label modulo 256, and join any two
distinct observed labels in the same output lane whose counts
sum to at least nine. This can add false edges, but never removes
either endpoint of a true internal edge. For the 256 successive
aligned input rows, the nominal middle-byte label before the final
addition runs once through all 256 values: each prior rotor's
middle quotient is a permutation, and each aligned quotient
translation is also a permutation. The true internal edges
therefore all appear. Edges crossing output-lane boundaries are
unneeded for the sixteen paths.

The implementation first enumerates paths in a lane with the
fewest observed edges. A completed path determines its unique
core \((u,v,w,t)\). Testing its sixteen input XORs determines each
possible common external input/output XOR pair. For every other
lane, equation~\eqref{eq:middlecore} then leaves only the
\(16^2\) choices of \((u,w)\) with the fixed common \((v,t)\).
Test their fifteen edges and require distinct \(u\)'s. This
finite completion includes the true paths, even when the graph
has extra edges; it imposes no path-count limit. Applying
Theorem~\ref{thm:forty} to every completed path set includes
the true raw rotor key.

For the numerical row test, invert a candidate \(K_8\) on each
of the 256 rows. For a proposed last residue \(b\), count
inverse images in the disjoint intervals
\([16v+b,16v+b+15]\). Keep \(b\) if every row has at least
\begin{equation}
                    17-\min(b,16-b)
\label{eq:carrytest}
\end{equation}
images in one such interval. The same \(b\) must work for all
rows. On the event, \(S\le\min(b,16-b)-1\), so
Lemma~\ref{lem:block} gives at least \(16-S\) images in the
true interval, which meets \eqref{eq:carrytest}.
All keys passing this exact necessary test are retained.

For the count, let \(a_d(s)\) be the number of \(d\)-dimensional
signed residue vectors of absolute norm \(s\). Then
\[
a_d(0)=1,\qquad
a_d(s)=\sum_{j=1}^{\min(d,s)}
2^j \binom dj\binom{s-1}{j-1}\quad(s>0).
\]
Choose the \(j\) nonzero coordinates, their signs, and a positive
composition of \(s\). For a fixed \(s\le7\), exactly \(15-2s\)
values of \(b\) obey \eqref{eq:outerevent}. Hence
\[
\sum_{s=0}^7 a_6(s)(15-2s)
=15+156+792+2628+6384+11820+16008+10836
=48639.
\]
Lemma~\ref{lem:uniform} divides this count by \(16^7\).
\end{proof}

The public screen samples 31 consecutive words at each of twelve
fixed input positions. Every candidate alignment has a complete
row in a 31-word strip, and on the event its two supported labels
sum to at least nine. Thus any one strip preserves its alignment.
The twelve strips only reject more IVs before the 4111-word
partial codebook is requested. That codebook contains all 256
rows for every alignment. The screen does not inspect a key or a
hidden state.

\section{One inward factorization}\label{sec:inward}
Write each state word \(s_i=b_i+256h_i\), \(0\le b_i,h_i<256\),
and let \(h_9=0\). Define the translated prefix
\[
                       C_t(p)=R_t(p)+256h_{t+1}.
\]
It obeys the same step at every stage:
\begin{equation}
 C_{t-1}(p)=\D_t(C_t(p)-256h_{t+1})-b_t.
\label{eq:pivot}
\end{equation}
Indeed, substitution leaves \(R_{t-1}(p)+256h_t\).
The pivot's output translation only shifts quotient labels:
\[
\left\lfloor\frac{C_t(p)}{16}\right\rfloor
=\left\lfloor\frac{R_t(p)}{16}\right\rfloor+16h_{t+1}
\pmod {4096}.
\]
It preserves the low-nibble block partition used above.

For each \(t=6,5,4,3\), define the inward event \(E_t\):
all low-nibble residues \(s_2,\ldots,s_6\) vanish except
\(s_t\bmod16=e_t\), where \(e_t\) is signed with
\(1\le|e_t|\le3\), and
\begin{equation}
           2|e_t|<b_7<16-2|e_t|,
           \qquad b_7=s_7\bmod16.
\label{eq:inner}
\end{equation}
Here \(s_1\) is handled by an unknown row alignment. For
each \(t\), the choices of \((e_t,b_7)\) number
\(2(11+7+3)=42\), so
\begin{equation}
                  q=\Pr(E_t)=42/16^6.
\label{eq:q}
\end{equation}
The four events are disjoint because their sole nonzero residue
is in a different state word.

Under any \(E_t\), Lemma~\ref{lem:block} exposes a translated
stage-seven quotient cycle. Trying the sixteen possible low
output residues and row alignments, the complete-cycle decoder
of Section~\ref{sec:decoder} includes \(K_7\) and its output
translation \(s_8\). The implementation first checks 256
partial rows, then asks for the full first-word codebook only
when a compatible numerical candidate and one next-stage
graph remain. Every test is a necessary condition on the
sufficient event. The full row test keeps every compatible
state residue; it does not retain just the best one.
The saved public context is
\begin{equation}
                 Y(p)=\D_7(\D_8(C(p))-s_8)=R_6(p)+s_7,
\label{eq:context}
\end{equation}
together with its IV, alignment, and admissible low residues
of \(s_7\).

For \(E_6\), inverting the recovered \(K_7\) exposes the
translated stage-six cycle as well. The exact cycle decoder
and full row test give seed pairs \((K_7,K_6)\).
Numerical suffixes obtained from one IV are pooled across all
saved contexts. The true suffix can therefore be extended using
a different IV for each inward event. All state offsets stay
with the IV from which they were obtained.

\subsection{The partition join}
Let \(q_j(x)=\E_j(16x)\gg4\) be the 4096-point quotient of
rotor \(j\), and let \(\Tr_a(x)=x+a\pmod {4096}\).
Suppose \(K_{t+1},\ldots,K_6\) are known and put \(d=6-t\),
so \(d\in\{1,2,3\}\). Under \(E_t\), the quotient walk of a
compatible context has the form
\begin{equation}
 W=\Tr_{\eta_d}q_{t+d}\Tr_{\eta_{d-1}}q_{t+d-1}
       \cdots \Tr_{\eta_1}q_{t+1}\Tr_{\eta_0}q_t,
\label{eq:walk}
\end{equation}
up to a cyclic input origin and direction. The \(\eta_j\)
are the state quotient translations carried by
\eqref{eq:pivot}; all are enumerated. The only nonaligned
low-nibble addition is the small \(e_t\), which creates the
two-label row edges. Every known quotient and translation
in \eqref{eq:walk} permutes the middle byte.

For \(0\le j\le d\), put
\(m_j(x)=q_{t+j}(x)\bmod256\) on \(0\le x<256\).
Lemma~\ref{lem:triangular} makes this the complete middle-byte
quotient, independent of the input top nibble.
Let \(\calP\) partition \(\{0,\ldots,255\}\) into its sixteen
consecutive 16-point blocks. For any 256-point permutation
\(f\), \(f(\calP)\) means the unordered set of image blocks.
Write \(m_j\) for the middle-byte quotient of the corresponding
\(q_j\). The exact join uses the identity
\begin{equation}
\left(\Tr_{-z}m_d^{-1}\Tr_{-h}W\Tr_a\right)(\calP)
=
\left(\Tr_{16H}m_{d-1}\Tr_{\theta_{d-2}}\cdots
       m_1\Tr_f\right)(\calP).
\label{eq:partition}
\end{equation}
All maps in this equation are reduced modulo 256.
For \(d=1\), the right side is simply \(\calP\).
Here \(h\in\{0,\ldots,255\}\) is the outer output
translation, \(z,a,f,H\in\{0,\ldots,15\}\), and the
\(\theta_j\in\{0,\ldots,255\}\) are inner translations
that exist only at greater depth. The left side has
\(256\cdot16\cdot16=65536\) choices; the right has
respectively \(1,256,65536\) choices for depths \(1,2,3\).

\begin{lemma}[Exact join]\label{lem:join}
On the true branch of any \(E_t\), one pair of entries in
\eqref{eq:partition} matches. The implemented join enumerates
every matching pair, including \(f=0\). For each match, its
middle-path and top-coordinate tests emit every compatible
numerical \(K_t\) and quotient translation.
\end{lemma}
\begin{proof}
Choose \(h,z,a\) to cancel the true outer translation and align
the input origin with a block boundary. Undo the known outer
\(m_d\). Decompose the next translation as \(z+16H\)
and the translation next to \(m_t\) as \(f+16F\).
Translations by multiples of sixteen preserve \(\calP\);
\(m_t\) preserves \(\calP\) by Lemma~\ref{lem:triangular}.
Moving those two factors inside the image of \(\calP\)
gives exactly the right side of \eqref{eq:partition};
the intervening 256-point translations are the
\(\theta_j\). This argument does not require \(f\ne0\).

The code encodes a partition by naming its sixteen blocks
in order of first occurrence among output positions, then
packs the 256 names into 128 bytes. Two codes are equal
exactly when the unordered partitions are equal. It sorts
the right family and compares full encodings, not hashes;
every matching pair is processed. Undoing its known factors
leaves the sixteen paths of an unknown middle map, with a
common output-lane translation. Theorem~\ref{thm:forty}
enumerates every matching numerical key.

It remains to restore the top quotient nibble. Write
\(q_j(256A+x)=256P_{j,x}(A)+m_j(x)\), with \(0\le A<16\).
Once the middle trajectory is fixed, the carries at all
known translations and the lane \(x\) of each known rotor
are fixed. On the sixteen inputs with one chosen middle
coordinate, the top trajectory is a composition of known
nibble permutations \(P_{j,x}\) and the remaining unknown
top-nibble additions. Split this composition in two.
Enumerate the at most 256 offset choices on one side and
the at most sixteen on the other; compare their complete
16-entry maps up to cyclic input and output translations.
For every equality, enumerate the sixteen input origins
and compare all 4096 entries of \eqref{eq:walk}.
The true offsets give an equality and pass that literal
comparison. No match-count assumption is used.
\end{proof}

The same join is called for \(d=1,2,3\), yielding \(K_5,K_4,K_3\).
After its quotient comparison, the code undoes the known
full-word rotors and tests every aligned row against the
candidate inner rotor's translated interval. On \(E_t\)
there are no low-nibble perturbations before \(t\), so
Lemma~\ref{lem:block} makes the true row test exact.
Every surviving offset and key is emitted.

\section{Initialization anchors and the first two rotors}
The partition join stops at \(K_3\). The remaining two rotors
are recovered from four rounds of initialization, without an
unbounded search over the other state words.
Let \(u_i^{(r)}\) be state word \(i\) after initialization round
\(r\), with \(u_i^{(0)}=\mathrm{IV}_i\).
Lemma~\ref{lem:uniform} gives, for \(i\ge3\),
\begin{equation}
 u_i^{(r-1)}
=u_i^{(r)}-\E_{i-1}(u_{i-1}^{(r)}).
\label{eq:cone}
\end{equation}
Known tail states and keys fill this inverse triangle.
In particular, the four known tail states give
\begin{equation}
 s_1=\mathrm{IV}_1+
       \sum_{r=1}^4\E_8(u_8^{(r)}).
\label{eq:s1}
\end{equation}

For \(t=2,3,4\), the same triangle computes
\(A_t=\E_t(s_t)\) from \(s_{t+1},\ldots,s_{t+4}\),
\(K_{t+1},\ldots,K_{t+3}\), and \(\mathrm{IV}_{t+4}\),
without consulting \(K_t\). To make the construction explicit,
set \(a_{0,j}=s_{t+1+j}\) for \(0\le j\le3\), then
\begin{equation}
 a_{r,j}=a_{r-1,j}-
 \E_{t+j}(a_{r-1,j-1}),
 \qquad 1\le r\le j\le3.
\label{eq:anchortriangle}
\end{equation}
Induction using \eqref{eq:cone} gives
\(a_{r,j}=u_{t+1+j}^{(4-r)}\). Starting from
\(x=\mathrm{IV}_{t+4}\), replace
\(x\leftarrow\D_{t+j}(a_{j,j}-x)\) for \(j=3,2,1\).
The result is the previous-round value feeding \(s_{t+1}\),
so \(A_t=s_{t+1}-x\).

Once \(K_3,s_4,\ldots,s_8\) are known on a true branch,
\eqref{eq:s1} and \(A_3\) give \(s_1\) and
\(s_3=\D_3(A_3)\). The triangle gives \(A_2\).
Undo the known rotors to obtain the complete public table
\begin{equation}
 R_2(p)=\E_2\bigl(\E_1(p+s_1)+s_2\bigr),
 \qquad A_2=\E_2(s_2).
\label{eq:pair}
\end{equation}
The final sufficient event also requires
\(s_2\bmod256\ne0\), with independent conditional
probability \(15/16\).

\begin{lemma}[Anchored pair retention]\label{lem:pair}
Given the correct \(R_2\) table, \(s_1,A_2\), and
\(s_2\bmod256\ne0\), the finite pair decoder emits the
true \((K_1,K_2,s_2)\). It then checks the complete
word table and the literal IV initialization.
\end{lemma}
\begin{proof}
Align the 256-point input high-byte rows by \(s_1\).
The first rotor maps each row to one high-byte class.
Adding a nonzero low byte of \(s_2\) splits it between two
adjacent classes, and the second rotor permutes those
classes. Thus every row has exactly two high-byte labels.
The resulting 256 edges trace the high-byte quotient
cycle of \(\E_2\), with the common split count giving
\(s_2\bmod256\).

The unique preimage of \(A_2\) in \eqref{eq:pair}
is the row where the first rotor output is zero.
It anchors the cycle origin and orientation. The
aligned row labels then give the first rotor's complete
high-byte quotient, while the anchored cycle gives
the second rotor's high-byte quotient up to one of
256 input origins. The decoder enumerates every such
origin. For each proposed quotient, it enumerates the
raw-key lifts and checks its complete 256-entry truth
table; the true raw key is among them by construction.

For clarity, the finite high-byte decoder uses the
autonomous \(Q\) of Lemma~\ref{lem:triangular} and two
top-coordinate rows \(P_b\). Put
\(p_0(b)=S_2(b\xor\beta_0)\) and
\(p_j(b)=S_2(p_{j-1}(b)\xor\beta_j)\) for \(j=1,2,3\).
Direct substitution gives
\[
 Q(b)=S_2(p_3(b)\xor\beta_0\xor\beta_1)
            \xor\beta_2\xor\beta_3.
\]
Its 65536 possible 16-entry tables are distinct.
The top rows are XOR translates of
\[
 J_P^{u,v,w,t}(x)
=S_1(S_3(S_1(S_3(S_1(x)\xor u)\xor v)\xor w)\xor t),
\]
with
\[
 (u,v,w,t)=
 (\gamma_1,\alpha_2\xor p_1(b),\gamma_3,
  \alpha_0\xor\alpha_1\xor p_3(b)).
\]
Exhaustively canonicalizing these 65536 nibble tables
under input/output XOR gives 65535 orbits: one two-core
collision and three self-stabilizing cores. The five
exceptional cores have distinct \((u,w)\) pairs.
The two rows \(b=0,1\) share \((u,w)\) but have different
\(v\), since \(p_1\) is a permutation. Thus at least
one row is nonexceptional and fixes its core and
external XORs. The raw schedule then leaves at most
sixteen low-nibble choices, all scanned. These finite
facts are reproduced by evaluating the printed
16-entry formulas over their complete parameter
ranges and sorting the resulting packed truth tables.

For the first rotor there are at most sixteen such
raw lifts; for the second there are at most
\(256\cdot16=4096\) origin/lift pairs. The decoder
checks each pair against all 65536 entries of
\eqref{eq:pair}, then checks the eight initialized
state words against the supplied IV. The true pair
passes every comparison.
\end{proof}

\section{True-key survival}\label{sec:bounds}
The attack has one induction. Acquire outer IVs
until the fixed budget has been examined; retain every outer
numerical key from Section~\ref{sec:outer}. For each such key,
acquire the fixed inward budget, save every stage-seven
context, and pool every seed pair. Apply the same partition
join at depths one, two, and three. At each leaf, use the
initialization anchors and anchored pair decoder, then
replay every complete key against every saved public
codebook. No score, candidate cap, or choice of a single
state branch occurs in this induction.

\begin{theorem}[True-key coverage]\label{thm:coverage}
For every fixed 256-bit master key, use \(m\) independent
uniform outer IVs and \(N\) independent uniform inward IVs
for each outer candidate. The probability that the final
list omits the true master key is at most
\begin{equation}
 \epsilon(m,N)
=\left(1-\frac{48639}{16^7}\right)^m
 +3\left(1-\frac{42}{16^6}\right)^N
 +\left(1-\frac{630}{16^7}\right)^N .
\label{eq:epsilon}
\end{equation}
At \(m=132000\), \(N=8000000\), this is
\begin{equation}
             \epsilon(m,N)
             =1.30716322171\cdot10^{-8}
             <1.4020\cdot10^{-8}.
\label{eq:number}
\end{equation}
The bound concerns inclusion of the true key in the
returned list.
\end{theorem}
\begin{proof}
By Theorem~\ref{thm:outer}, one outer event retains
the true \(K_8\). On that branch, \(E_6\) retains the
true seed \((K_7,K_6)\), and \(E_5,E_4,E_3\) extend
the true numerical suffix in turn. The proof of
Lemma~\ref{lem:join} includes zero translations, so
there is no further state condition at \(E_4\).
For \(E_3\), Lemma~\ref{lem:pair} adds
\(s_2\bmod256\ne0\). Given \(E_3\), its upper
four bits of the low byte are unrestricted and
uniform by Lemma~\ref{lem:uniform}; this adds the
factor \(15/16\), giving \(42\cdot15/16^7
=630/16^7\). The correct branch passes literal
initialization and all public-ciphertext replay.

The absence probabilities for the outer event and
the four inward event families are respectively
\((1-p_{\rm out})^m\),
\((1-q)^N\) three times, and
\((1-630/16^7)^N\). Union bound them.
Independence among the different families is not
assumed. False branches can increase work but cannot
remove the true branch because every compatible
candidate is processed.
\end{proof}
\[
\boxed{\Pr\{\text{true master key is returned}\}
       \ge 1-1.30716322171\cdot10^{-8}.}
\]

\begin{thebibliography}{9}\small
\bibitem{separ} A. Vahi and S. Jafarali Jassbi.
SEPAR: A New Lightweight Hybrid Encryption Algorithm with a Novel
Design Approach for IoT. \emph{Wireless Personal Communications}
114, 2283--2314, 2020.
\href{https://doi.org/10.1007/s11277-020-07476-y}
{doi:10.1007/s11277-020-07476-y}.
\end{thebibliography}
\end{document}
